\subsection{Dynkin diagrams}\label{sga3:XXI:7.4}

\begin{definition}\label{sga3:XXI:7.4.1}
A \emph{Dynkin-diagram structure} on a finite set \(\Delta\)
(the word ``scheme'' being banished for obvious reasons) consists
of a set of pairs of distinct elements of \(\Delta\), called
\emph{linked pairs}, together with a map
\[
 \Delta\longrightarrow\{1,2,3\}.
\]
The notion of isomorphism between such structures is evident.
\end{definition}

\begin{definition}\label{sga3:XXI:7.4.2}
Let \(\cR\) be a root datum and let \(\Delta\) be a system of
simple roots.  The \emph{Dynkin diagram} of \(\cR\) relative to
\(\Delta\) is the set \(\Delta\), in which two simple roots are
linked if and only if they are nonorthogonal, and each root is
assigned its length.
\end{definition}

\begin{proposition}\label{sga3:XXI:7.4.3}
The Dynkin diagram and the Cartan matrix determine one another
uniquely.
\end{proposition}

\begin{proof}
The equivalence
\[
 \alpha\text{ and }\beta\text{ are not linked}
 \quad\Longleftrightarrow\quad
 a_{\alpha,\beta}=0
\]
and the relation
\[
 \operatorname{long}(\alpha)a_{\alpha,\beta}
 =
 \operatorname{long}(\beta)a_{\beta,\alpha},
\]
with \(\inf\operatorname{long}(\alpha)=1\) in each connected
component of the diagram, determine the \(a_{\alpha,\beta}\) from
the links and lengths, and conversely.  The details are left to the
reader.
\end{proof}

\begin{corollary}\label{sga3:XXI:7.4.4}
A reduced simply connected or adjoint root datum is determined by
its Dynkin diagram.
\end{corollary}

\begin{corollary}\label{sga3:XXI:7.4.5}
Let \(\cR\) be a reduced simply connected or adjoint root datum and
let \(\Delta\) be a system of simple roots.  The group
\(E_\Delta(\cR)\) is identified with the group of automorphisms of
the Dynkin diagram of \(\cR\) relative to \(\Delta\); that is, with
the group of permutations of \(\Delta\) preserving lengths and
links.
\end{corollary}

\begin{remark}\label{sga3:XXI:7.4.6}
The usual method\footnote{See Bourbaki, \emph{Groupes et algèbres
de Lie}, Chapter~VI, no.~4.2, or the Sophus Lie Seminar.}
\footnote{Editors' note: for a slightly different proof, see also
[De80].}
classifies the connected Dynkin diagrams and shows that each
actually corresponds to an irreducible reduced simply connected
root datum.  One obtains the familiar types:
\begin{center}
\begin{tikzpicture}[
  x=.86cm,y=.78cm,
  dyn/.style={circle,fill=black,inner sep=1.35pt},
  wt/.style={font=\small,above=2pt},
  typ/.style={font=\normalsize,anchor=west}
]
\foreach \x in {0,1,2,4,5,6} {
  \node[dyn] (A\x) at (\x,0) {};
  \node[wt] at (A\x) {$1$};
}
\draw (A0)--(A1)--(A2);
\draw (A4)--(A5)--(A6);
\node at (3,0) {$\cdots$};
\node[typ] at (7.3,0) {$A_n,\qquad n\geq1.$};

\foreach \x in {0,1,2,4,5} {
  \node[dyn] (B\x) at (\x,-1.55) {};
  \node[wt] at (B\x) {$2$};
}
\node[dyn] (B6) at (6,-1.55) {};
\node[wt] at (B6) {$1$};
\draw (B0)--(B1)--(B2);
\draw (B4)--(B5)--(B6);
\node at (3,-1.55) {$\cdots$};
\node[typ] at (7.3,-1.55) {$B_n,\qquad n\geq2.$};

\foreach \x in {0,1,2,4,5} {
  \node[dyn] (C\x) at (\x,-3.10) {};
  \node[wt] at (C\x) {$1$};
}
\node[dyn] (C6) at (6,-3.10) {};
\node[wt] at (C6) {$2$};
\draw (C0)--(C1)--(C2);
\draw (C4)--(C5)--(C6);
\node at (3,-3.10) {$\cdots$};
\node[typ] at (7.3,-3.10) {$C_n,\qquad n\geq3.$};

\foreach \x in {0,1,2,4} {
  \node[dyn] (D\x) at (\x,-4.65) {};
  \node[wt] at (D\x) {$1$};
}
\node[dyn] (Du) at (5,-4.25) {};
\node[wt] at (Du) {$1$};
\node[dyn] (Dl) at (5,-5.05) {};
\node[wt] at (Dl) {$1$};
\draw (D0)--(D1)--(D2);
\draw (D4)--(Du);
\draw (D4)--(Dl);
\node at (3,-4.65) {$\cdots$};
\node[typ] at (7.3,-4.65) {$D_n,\qquad n\geq4.$};

\foreach \x in {0,1,3,4,5} {
  \node[dyn] (E\x) at (\x,-6.35) {};
  \node[wt] at (E\x) {$1$};
}
\node[dyn] (Ev) at (3,-7.15) {};
\node[right=2pt,font=\small] at (Ev) {$1$};
\draw (E0)--(E1);
\draw (E3)--(E4)--(E5);
\draw (E3)--(Ev);
\node at (2,-6.35) {$\cdots$};
\node[typ] at (7.3,-6.65) {$E_n,\qquad n=6,7,8.$};

\foreach \x/\w in {0/1,1/1,2/2,3/2} {
  \node[dyn] (F\x) at (\x,-8.15) {};
  \node[wt] at (F\x) {$\w$};
}
\draw (F0)--(F1)--(F2)--(F3);
\node[typ] at (7.3,-8.15) {$F_4.$};

\node[dyn] (G0) at (2,-9.70) {};
\node[wt] at (G0) {$1$};
\node[dyn] (G1) at (3,-9.70) {};
\node[wt] at (G1) {$3$};
\draw (G0)--(G1);
\node[typ] at (7.3,-9.70) {$G_2.$};
\end{tikzpicture}

\end{center}
By \ref{sga3:XXI:7.4.5}, the corresponding group
\(E_\Delta(\cR)\) is obtained at once:
\[
\begin{array}{ll}
E_\Delta(\cR)=\{e\},
 & A_1,\ B_n,\ C_n,\ E_7,\ E_8,\ F_4,\ G_2,\\[2pt]
E_\Delta(\cR)=\mathbb Z/2\mathbb Z,
 & A_n\ (n\geq2),\ D_n\ (n\geq5),\ E_6,\\[2pt]
E_\Delta(\cR)=\mathfrak S_3,
 & D_4.
\end{array}
\]
\end{remark}
